\documentclass[12pt,a4paper]{article}

\usepackage{cmap}
\usepackage[T2A]{fontenc}
\usepackage[utf8]{inputenc}
\usepackage[russian]{babel}
\usepackage{amsmath,amssymb}
\usepackage[a4paper,margin=2cm]{geometry}

\newcommand{\R}{\mathbb{R}}
\newcommand{\CC}{\mathbb{C}}
\newcommand{\N}{\mathbb{N}}
\newcommand{\dd}{\,d}

\setlength{\parindent}{0pt}
\setlength{\parskip}{0.45em}
\sloppy

\begin{document}

\begin{center}
{\Large\bfseries Билет 13}
\end{center}

\noindent\fbox{%
  \begin{minipage}{\dimexpr\linewidth-2\fboxsep-2\fboxrule\relax}
  \normalfont
Степенные ряды с действительными членами. Почленное интегрирование
степенного ряда. Бесконечная дифференцируемость суммы степенного ряда на
интервале сходимости. Единственность представления функции степенным рядом.
Достаточные условия разложимости бесконечно дифференцируемой функции в
степенной ряд. Ряд Тейлора. Формула Тейлора с остаточным членом в
интегральной форме. Пример бесконечно дифференцируемой функции, не
разлагающейся в степенной ряд. Разложение в ряд Тейлора основных
элементарных функций:
\(e^x\), \(\cos x\), \(\sin x\), \(\ln(1+x)\), \((1+x)^\alpha\). Разложение в степенной
ряд функции \(e^z\) комплексного переменного \(z\).
  \end{minipage}%
}

\section*{Краткий план}

Сначала степенной ряд с действительными коэффициентами рассматривается как
частный случай комплексного степенного ряда, поэтому для него есть радиус и
интервал сходимости. Внутри интервала сходимости ряд сходится равномерно на
любом вложенном отрезке; отсюда следуют почленное интегрирование,
почленное дифференцирование любого порядка и единственность коэффициентов.
Затем вводится ряд Тейлора и понятие регулярности: функция разлагается в
степенной ряд тогда, когда остаток формулы Тейлора стремится к нулю. Для
контроля остатка используется формула Тейлора в интегральной форме и
достаточное условие ограниченности всех производных. В конце записываются
стандартные разложения элементарных функций и пример гладкой, но не
аналитической функции.

\section*{Определения}

\textbf{Вещественный степенной ряд.}
Рядом с действительными членами с центром \(x_0\) называется ряд
\[
  \sum_{k=0}^{\infty} a_k(x-x_0)^k,
  \qquad a_k,x,x_0\in\R.
\]
Его радиус сходимости \(R\in[0,+\infty]\) можно находить по формуле
Коши--Адамара
\[
  \frac1R=\limsup_{k\to\infty}\sqrt[k]{|a_k|},
  \qquad \frac10=+\infty,\quad \frac1{+\infty}=0.
\]
Интервал \((x_0-R,x_0+R)\) называется интервалом сходимости. Внутри него
ряд сходится абсолютно, вне замкнутого интервала расходится; в граничных
точках вопрос решается отдельно.

\textbf{Сумма степенного ряда.}
Если \(R>0\), то на интервале сходимости определена функция
\[
  f(x)=\sum_{k=0}^{\infty}a_k(x-x_0)^k.
\]

\textbf{Ряд Тейлора.}
Если функция \(f\) бесконечно дифференцируема в точке \(x_0\), то ряд
\[
  \sum_{k=0}^{\infty}\frac{f^{(k)}(x_0)}{k!}(x-x_0)^k
\]
называется рядом Тейлора функции \(f\) в точке \(x_0\). При \(x_0=0\) это
ряд Маклорена.

\textbf{Регулярность в точке.}
Функция \(f\) называется регулярной в точке \(x_0\), если она бесконечно
дифференцируема в этой точке и ее ряд Тейлора сходится к \(f\) в некоторой
окрестности \(x_0\):
\[
  \exists\delta>0:\quad
  f(x)=\sum_{k=0}^{\infty}\frac{f^{(k)}(x_0)}{k!}(x-x_0)^k
  \quad (x\in U_\delta(x_0)).
\]

\textbf{Остаточный член формулы Тейлора.}
Для \(n\) раз дифференцируемой функции
\[
  S_n(x)=\sum_{k=0}^{n}\frac{f^{(k)}(x_0)}{k!}(x-x_0)^k,\qquad
  r_n(x)=f(x)-S_n(x).
\]
Функция регулярна в \(x_0\) тогда и только тогда, когда для некоторого
\(\delta>0\)
\[
  x\in U_\delta(x_0)\quad\Longrightarrow\quad r_n(x)\to0.
\]

\section*{Теоремы}

\textbf{1. Радиус не меняется при формальном дифференцировании и
интегрировании.}
Если степенной ряд \(\sum c_kz^k\) имеет радиус сходимости \(R\), то ряды
\[
  \sum_{k=1}^{\infty}kc_kz^{k-1},
  \qquad
  \sum_{k=0}^{\infty}\frac{c_k}{k+1}z^{k+1}
\]
имеют тот же радиус сходимости \(R\).

\textbf{2. Почленное интегрирование, дифференцирование и единственность.}
Пусть
\[
  f(x)=\sum_{k=0}^{\infty}a_k(x-x_0)^k
\]
имеет радиус сходимости \(R>0\). Тогда для всякого
\(x\in(x_0-R,x_0+R)\)
\[
  \int_{x_0}^{x} f(t)\dd t
  =
  \sum_{k=0}^{\infty}\frac{a_k}{k+1}(x-x_0)^{k+1}.
\]
Кроме того, \(f\) имеет производные всех порядков внутри интервала сходимости,
и
\[
  f^{(n)}(x)
  =
  \sum_{k=n}^{\infty}
  a_k k(k-1)\cdots(k-n+1)(x-x_0)^{k-n}.
\]
В частности,
\[
  a_n=\frac{f^{(n)}(x_0)}{n!},
\]
поэтому представление функции степенным рядом с данным центром единственно.

\textbf{3. Формула Тейлора с интегральным остатком.}
Если \(f\) имеет непрерывные производные до порядка \(n+1\) включительно в
окрестности \(U_\delta(x_0)\), то
\[
  f(x)=\sum_{k=0}^{n}\frac{f^{(k)}(x_0)}{k!}(x-x_0)^k+r_n(x),
\]
где
\[
  r_n(x)=\frac1{n!}\int_{x_0}^{x}(x-t)^n f^{(n+1)}(t)\dd t,
  \qquad x\in U_\delta(x_0).
\]

\textbf{4. Достаточное условие разложимости в степенной ряд.}
Пусть существует \(\delta>0\), такая что \(f\in C^\infty(U_\delta(x_0))\), и
\[
  \exists M>0:\quad
  |f^{(n)}(x)|\le M
  \quad(n\in\N,\ x\in U_\delta(x_0)).
\]
Тогда \(f\) регулярна в \(x_0\), то есть
\[
  f(x)=\sum_{k=0}^{\infty}\frac{f^{(k)}(x_0)}{k!}(x-x_0)^k
  \quad (x\in U_\delta(x_0)).
\]

\textbf{5. Основные разложения.}
\[
  e^x=\sum_{k=0}^{\infty}\frac{x^k}{k!},
  \qquad x\in\R,
\]
\[
  \cos x=\sum_{k=0}^{\infty}(-1)^k\frac{x^{2k}}{(2k)!},
  \qquad
  \sin x=\sum_{k=0}^{\infty}(-1)^k\frac{x^{2k+1}}{(2k+1)!},
  \qquad x\in\R,
\]
\[
  \ln(1+x)=\sum_{n=1}^{\infty}(-1)^{n-1}\frac{x^n}{n},
  \qquad -1<x\le1,
\]
\[
  (1+x)^\alpha=\sum_{k=0}^{\infty}C_\alpha^k x^k,
  \qquad -1<x<1,
\]
где
\[
  C_\alpha^0=1,\qquad
  C_\alpha^k=\frac{\alpha(\alpha-1)\cdots(\alpha-k+1)}{k!}.
\]
Если \(\alpha\in\N\cup\{0\}\), то ряд конечен и равен биному Ньютона при всех
\(x\in\R\). Для комплексной экспоненты
\[
  e^z=\sum_{k=0}^{\infty}\frac{z^k}{k!},
  \qquad z\in\CC.
\]

\section*{Доказательства}

\textbf{Радиус сходимости производного и интегрального рядов.}
Для ряда \(\sum c_kz^k\) по формуле Коши--Адамара
\[
  \frac1{R_1}
  =
  \limsup_{k\to\infty}\sqrt[k]{|kc_k|}
  =
  \lim_{k\to\infty}\sqrt[k]{k}\,
  \limsup_{k\to\infty}\sqrt[k]{|c_k|}
  =
  \frac1R,
\]
поскольку \(\sqrt[k]{k}\to1\). Ряды
\(\sum_{k=1}^{\infty}kc_kz^k\) и
\(\sum_{k=1}^{\infty}kc_kz^{k-1}\) при \(z\ne0\) отличаются умножением суммы
частичных сумм на \(z\), а при \(z=0\) оба сходятся. Значит радиус
производного ряда равен \(R\). Интегральный ряд проверяется так же, потому
что исходный ряд получается из него почленным дифференцированием.

\textbf{Почленное интегрирование и дифференцирование.}
Пусть \(x\) лежит внутри интервала сходимости. Выберем \(r<R\) так, что
\(x\in[x_0-r,x_0+r]\). Степенной ряд равномерно сходится на этом отрезке.
По теореме о почленном интегрировании равномерно сходящегося ряда
\[
  \int_{x_0}^{x}f(t)\dd t
  =
  \sum_{k=0}^{\infty}a_k\int_{x_0}^{x}(t-x_0)^k\dd t
  =
  \sum_{k=0}^{\infty}\frac{a_k}{k+1}(x-x_0)^{k+1}.
\]
Производный степенной ряд имеет тот же радиус \(R\), поэтому он также
равномерно сходится на \([x_0-r,x_0+r]\). По теореме о почленном
дифференцировании функционального ряда
\[
  f'(x)=\sum_{k=1}^{\infty}ka_k(x-x_0)^{k-1}.
\]
Повторяя рассуждение, получаем производные любого порядка. При \(x=x_0\)
в формуле для \(f^{(n)}\) все слагаемые исчезают, кроме \(k=n\), поэтому
\(f^{(n)}(x_0)=n!a_n\). Отсюда следует единственность коэффициентов.

\textbf{Интегральный остаток формулы Тейлора.}
При \(n=0\)
\[
  r_0(x)=f(x)-f(x_0)=\int_{x_0}^{x}f'(t)\dd t,
\]
то есть формула верна. Пусть она доказана для \(n=s-1\):
\[
  r_{s-1}(x)=\frac1{(s-1)!}\int_{x_0}^{x}(x-t)^{s-1}f^{(s)}(t)\dd t.
\]
Интегрируя по частям, получаем
\[
  r_{s-1}(x)
  =
  \frac{f^{(s)}(x_0)}{s!}(x-x_0)^s
  +
  \frac1{s!}\int_{x_0}^{x}(x-t)^sf^{(s+1)}(t)\dd t.
\]
После вычитания \(s\)-го члена многочлена Тейлора получается формула для
\(r_s(x)\). Индукция завершает доказательство.

\textbf{Достаточное условие регулярности.}
По формуле Тейлора с остатком в форме Лагранжа для
\(x\in U_\delta(x_0)\)
\[
  r_n(x)=\frac{f^{(n+1)}(\xi)}{(n+1)!}(x-x_0)^{n+1},
\]
где \(\xi\) лежит между \(x\) и \(x_0\). Поэтому
\[
  |r_n(x)|
  \le
  M\frac{\delta^{n+1}}{(n+1)!}\to0.
\]
Значит ряд Тейлора сходится к \(f\) в \(U_\delta(x_0)\).

\textbf{Пример бесконечно дифференцируемой функции, не разлагающейся в
степенной ряд.}
Положим
\[
  f(x)=
  \begin{cases}
    e^{-1/x^2},& x\ne0,\\
    0,& x=0.
  \end{cases}
\]
Для любого \(m\in\N\)
\[
  \lim_{x\to0}x^{-m}e^{-1/x^2}=0.
\]
Индукцией по \(n\) получается
\[
  f^{(n)}(x)=
  \begin{cases}
    P_{3n}(1/x)e^{-1/x^2},& x\ne0,\\
    0,& x=0,
  \end{cases}
\]
где \(P_{3n}\) -- некоторый многочлен. Следовательно,
\(f^{(n)}(0)=0\) для всех \(n\), и ряд Тейлора в нуле тождественно равен
нулю. Но \(f(x)>0\) при \(x\ne0\), поэтому в любой окрестности нуля сумма
ряда Тейлора не совпадает с функцией.

\textbf{Разложения элементарных функций.}
Для \(e^x,\sin x,\cos x\) все производные на каждом отрезке ограничены.
По достаточному условию регулярности эти функции равны своим рядам
Маклорена; коэффициенты вычисляются по значениям производных в нуле.

Для логарифма сначала при \(|x|<1\)
\[
  \frac1{1+x}=\sum_{k=0}^{\infty}(-1)^kx^k.
\]
Интегрируя этот степенной ряд от \(0\) до \(x\), получаем
\[
  \ln(1+x)=\sum_{k=0}^{\infty}(-1)^k\frac{x^{k+1}}{k+1}
  =
  \sum_{n=1}^{\infty}(-1)^{n-1}\frac{x^n}{n}.
\]
При \(x=1\) равенство сохраняется по второй теореме Абеля и непрерывности
суммы на \([0,1]\).

Для \((1+x)^\alpha\) применяем интегральную форму остатка к формуле
Маклорена. При \(-1<x<1\)
\[
  f^{(k)}(x)=\alpha(\alpha-1)\cdots(\alpha-k+1)(1+x)^{\alpha-k},
\]
и остаток имеет вид
\[
  r_n(x)=
  \lambda_n
  \int_0^1
  \left(\frac{1-\tau}{1+\tau x}\right)^n
  (1+\tau x)^{\alpha-1}\dd\tau,
\]
где
\[
  \lambda_n=
  \frac{\alpha(\alpha-1)\cdots(\alpha-n)}{n!}x^{n+1}.
\]
Так как
\[
  \frac{\lambda_{n+1}}{\lambda_n}
  =
  \frac{\alpha-n-1}{n+1}x\to -x,\qquad |x|<1,
\]
то \(\lambda_n\to0\), а интеграл ограничен по \(n\). Следовательно,
\(r_n(x)\to0\), и биномиальный ряд сходится к \((1+x)^\alpha\).

Для комплексной экспоненты ряд \(\sum z^k/k!\) имеет бесконечный радиус
сходимости. Обозначим его сумму через \(F(z)\). По произведению абсолютно
сходящихся рядов
\[
  F(z_1)F(z_2)=F(z_1+z_2).
\]
При вещественном \(x\) уже доказано \(F(x)=e^x\), а
\[
  F(iy)=
  \sum_{n=0}^{\infty}(-1)^n\frac{y^{2n}}{(2n)!}
  +i\sum_{n=0}^{\infty}(-1)^n\frac{y^{2n+1}}{(2n+1)!}
  =
  \cos y+i\sin y.
\]
Если \(z=x+iy\), то
\[
  F(z)=F(x)F(iy)=e^x(\cos y+i\sin y)=e^z.
\]

\section*{Финальная устная версия}

Вещественный степенной ряд имеет вид
\[
  \sum_{k=0}^{\infty}a_k(x-x_0)^k.
\]
Его радиус сходимости находится по Коши--Адамару, а интервал сходимости
равен \((x_0-R,x_0+R)\). Внутри этого интервала ряд равномерно сходится на
каждом вложенном отрезке. Поэтому его можно почленно интегрировать:
\[
  \int_{x_0}^{x}f(t)\dd t
  =
  \sum_{k=0}^{\infty}\frac{a_k}{k+1}(x-x_0)^{k+1},
\]
и почленно дифференцировать сколько угодно раз:
\[
  f^{(n)}(x)
  =
  \sum_{k=n}^{\infty}
  a_k k(k-1)\cdots(k-n+1)(x-x_0)^{k-n}.
\]
Подставляя \(x=x_0\), получаем \(a_n=f^{(n)}(x_0)/n!\), значит
представление степенным рядом единственно.

Ряд Тейлора бесконечно дифференцируемой функции в точке \(x_0\) равен
\[
  \sum_{k=0}^{\infty}\frac{f^{(k)}(x_0)}{k!}(x-x_0)^k.
\]
Функция разлагается в этот ряд в окрестности \(x_0\) тогда и только тогда,
когда остатки формулы Тейлора \(r_n(x)\) стремятся к нулю. Удобная формула
для остатка:
\[
  r_n(x)=\frac1{n!}\int_{x_0}^{x}(x-t)^nf^{(n+1)}(t)\dd t.
\]
Она доказывается индукцией по \(n\) с интегрированием по частям. Достаточное
условие регулярности: если все производные \(f^{(n)}\) в некоторой
окрестности \(x_0\) равномерно ограничены одним числом \(M\), то по формуле
Лагранжа
\[
  |r_n(x)|\le M\frac{\delta^{n+1}}{(n+1)!}\to0,
\]
и \(f\) равна своему ряду Тейлора.

Пример гладкой, но не разлагающейся функции:
\[
  f(x)=e^{-1/x^2}\ (x\ne0),\qquad f(0)=0.
\]
Все ее производные в нуле равны нулю, поэтому ряд Тейлора равен нулю, но сама
функция положительна при \(x\ne0\).

Основные разложения:
\[
  e^x=\sum_{k=0}^{\infty}\frac{x^k}{k!},\quad
  \cos x=\sum_{k=0}^{\infty}(-1)^k\frac{x^{2k}}{(2k)!},\quad
  \sin x=\sum_{k=0}^{\infty}(-1)^k\frac{x^{2k+1}}{(2k+1)!},
\]
\[
  \ln(1+x)=\sum_{n=1}^{\infty}(-1)^{n-1}\frac{x^n}{n},\quad -1<x\le1,
\]
\[
  (1+x)^\alpha=\sum_{k=0}^{\infty}C_\alpha^kx^k,\quad
  C_\alpha^k=\frac{\alpha(\alpha-1)\cdots(\alpha-k+1)}{k!},\quad -1<x<1,
\]
\[
  e^z=\sum_{k=0}^{\infty}\frac{z^k}{k!},\qquad z\in\CC.
\]

\section*{Источники}

Иванов Г.Е. \emph{Лекции по математическому анализу. Часть 1}:
\texttt{IvGE\_1.pdf}, стр. 304--321.

Основные роли источника: стр. 304--309 -- радиус и интервал сходимости,
равномерная сходимость степенного ряда, неизменность радиуса при формальном
дифференцировании и интегрировании; стр. 309--311 -- почленное
интегрирование, бесконечная дифференцируемость суммы и единственность
коэффициентов; стр. 311--313 -- ряд Тейлора, регулярность, пример
\(e^{-1/x^2}\), достаточное условие регулярности; стр. 314--316 -- ряды
Маклорена для \(e^x\), \(\cos x\), \(\sin x\) и разложение \(e^z\);
стр. 317--321 -- интегральный остаток формулы Тейлора, биномиальный и
логарифмический ряды.

Дополнительно был просмотрен \texttt{IvGE\_2.pdf}; прямого материала по этому
билету там не найдено, найденные совпадения были поздними ссылками или
нерелевантными темами.

\end{document}
